\section{Action, gravity, and thermal information}
\label{sec:gravity}
\subsection{Dodecaphasic period and circular realization of action}
The length in Section~\ref{g26:sec:rigidez-radial-en} fixes, within
the same realization, the compatible clock
\[
 t_0=\frac{\pi_{\HMT}L}{54c_{\rm int}}.
\]
Its circular representation then satisfies \(R_{108}=L\). This
substitution expresses the constructed length chronologically and retains
the longitudinal reference, action and velocity of that realization.
The calendar already constructed determines $T_{108}=108t_0$.
In its circular realization,
\[
\omega_{108}=\frac{2\pi}{108t_0},\qquad
R_{108}=\frac{c_{\rm int}}{\omega_{108}}
=\frac{54}{\pi}\ell_0,\qquad \ell_0=c_{\rm int}t_0.
\]
The unit $\ell_0$ and the time and velocity charts remain
explicit; the SI chart does not select the cycle.
For each oriented action section $a\in\{\mathrm{pre},\mathrm{ret}\}$,
\[
E_{108,a}=\hbar_a\omega_{108},\qquad
m_{108,a}=\frac{E_{108,a}}{c_{\rm int}^2}.
\]
The reduced and full lengths are, respectively,
\[
\frac{\hbar_a}{m_{108,a}c_{\rm int}}=R_{108},
\qquad
\frac{h_a}{m_{108,a}c_{\rm int}}=2\pi R_{108}.
\]
The same action and duration enter both readings
\cite{autodual}.

\begin{proposition}[Radius-coincidence selector]
In the dimensional collection
\[
G_a(\vartheta)=\vartheta\frac{c_{\rm int}^5t_0^2}{\hbar_a},
\qquad \vartheta>0,
\]
the circular-realization condition
$G_a(\vartheta)m_{108,a}/c_{\rm int}^2=R_{108}$
has the unique solution
\[
\vartheta^\circ_{108}=(54/\pi)^2.
\]
\end{proposition}
\begin{proof}
The first radius is
$r_g(\vartheta)=\vartheta(\pi/54)\ell_0$.
Comparing it with $(54/\pi)\ell_0$, with $\ell_0>0$, reduces the
equality to a strictly increasing linear equation.
\end{proof}
The coincidence condition is an explicit part of this realization;
the theorem solves for its coefficient without using a metrological value of $G$.
The radius used here is $Gm/c^2$, so the factor two from
the convention $2Gm/c^2$ is not inadvertently transferred.

\begin{proposition}[Reciprocal action and gravity, preserved area]
In the two preceding charts,
\[
\frac{G_{\rm pre}}{G_{\rm ret}}=R_{\rm act}^{-1},\qquad
\hbar_aG_a=\vartheta^\circ_{108}c_{\rm int}^5t_0^2,
\qquad
\ell_P^2=\frac{\hbar_aG_a}{c_{\rm int}^3}
=\vartheta^\circ_{108}\ell_0^2.
\]
\end{proposition}
\begin{proof}
The definition of $G_a$ contains $\hbar_a^{-1}$, and the remaining factors
are common to both charts. Taking the quotient and product gives
the three equalities.
\end{proof}
This reciprocity explains why orientation information may
vary while area remains unchanged. A precise relation is thus available
between the action ellipse, Planck scale, and area operator
of section~\ref{sec:area}.

\subsection{Thermal normalization and energy per cycle}
Fix the same oriented section \(a\) for energy and
temperature. The corpus's thermal relation gives
\begin{equation}
k_B\Theta_{{\rm clk},a}\log3
=\frac{h_a}{108t_0}
=\frac{\hbar_a}{t_P}=E_{P,a},
\qquad t_P=\frac{54}{\pi}t_0.
\label{eq:thermalclock}
\end{equation}
The factor $\log3$ measures the information in the equiprobable triple;
clock energy incorporates the action already generated.
The identity preserves the product $k_B\Theta_{\rm clk}$ as an energy
section and makes the thermal chart used explicit \cite{termica}.

At this scale, the radiative realization provides
\[
n_\gamma\ell_P^3=\frac{2\zeta(3)}{\pi^2\log^3 3},\qquad
\frac{u_\gamma\ell_P^3}{E_P}
=\frac{\pi^2}{15\log^4 3},\qquad
\frac{s_\gamma\ell_P^3}{k_B}
=\frac{4\pi^2}{45\log^3 3}.
\]
Two distinct invariants occur here: $\log3$ preserves TRIT
information and $\zeta(3)$ the bosonic occupation moment.
Their occurrence in the same expression relates the thermal constants
to the spectral moments of bosonic occupation, without identifying
these quantities with one another.

Section~\ref{sec:ii-kb-aut-desarrollo} develops the Boltzmann
transducer, its action on tritic information, and the change of
thermal coordinates. Equation~\eqref{eq:thermalclock} fixes the
energy--temperature product used by that construction. The
selection of an individual coordinate is distinguished there from the
identity for the product. Within this same development, the preceding radiative formulas
preserve the occupation and entropy content
of the bosonic realization.

\subsection{Barbero--Immirzi in the Cartan--Holst realization}
Once $\gamma_{\HMT}$ has been constructed, the source realizes it as a
coefficient of the Holst action. In the declared geometric chart,
let $e^I$ be a nondegenerate co-tetrad, $\omega^{IJ}$ a connection,
$\Sigma^{IJ}=e^I\wedge e^J$, and
$F^{IJ}=\dd\omega^{IJ}+\omega^I{}_K\wedge\omega^{KJ}$.
With $\kappa=8\pi G/c^4$, the specialized action is
\begin{equation}
\begin{aligned}
S[e,\omega,\Psi]
={}&\frac1{2\kappa c}\int
\Sigma_{IJ}\wedge(\star F^{IJ}+\gamma_{\HMT}^{-1}F^{IJ})\\
&-\frac{\Lambda}{\kappa c}\int\operatorname{vol}_e
+S_m[e,\omega,\Psi].
\end{aligned}
\label{eq:holst}
\end{equation}
The current convention is
$\delta_\omega S_m=-\tfrac12\int\delta\omega^{IJ}\wedge\sigma_{IJ}$.
For this variation, the connection is fixed at the boundary,
$\delta\omega^{IJ}|_{\partial M}=0$, or variations compactly
supported in the interior are used. These sufficient conditions annihilate
the boundary term resulting from integrating
$\Sigma_{IJ}\wedge(\star+\gamma_{\HMT}^{-1}I)
D_\omega\delta\omega^{IJ}$.
Variation of \eqref{eq:holst} under that condition gives
\[
D_\omega[(\star+\gamma_{\HMT}^{-1}I)\Sigma^{IJ}]
=\kappa c\,\sigma^{IJ}.
\]
This follows by using $\delta F=D_\omega\delta\omega$, integrating by parts,
and setting the coefficient of the arbitrary variation to zero
\cite{holst}. Boundary orientation and current
normalization are part of this formula.

In Lorentzian signature, $\star^2=-I$ on bivectors. The identity
\[
(\star+\gamma^{-1}I)^{-1}
=\frac{\gamma^2}{1+\gamma^2}(-\star+\gamma^{-1}I)
\]
is checked by multiplication: the cross terms
cancel, leaving $(1+\gamma^{-2})I$ before normalization.
Thus the parameter already constructed enters the torsional response,
not merely a table value.
The identities
$D_\omega T^I=F^I{}_J\wedge e^J$ and $D_\omega F^{IJ}=0$
keep torsion, curvature, and balance connected.

The geometric realization preserves route memory and boundary soldering.
The cosmological constant \(\Lambda\) in \eqref{eq:holst} is the
coefficient declared in the action; it is not identified with the vector
representation \(\Lambda(R)\). The following identities retain the
holonomy, smoothness, and nondegeneracy conditions of the realization.
% Article VII, additive recovery, revision 04, 10-09-2026.
% RESULTADO_RECUPERADO: S15, 11c_gravedad_holonomia_rev6.tex:16--126.
% The canonical cocycle c27=(1,18) is preserved as an antecedent of the route
% computation; this fragment proves its composition and representation.
% Finite translational memory and the cosmological amplitude s0 have
% different types. The cocycle does not by itself evaluate a local density.

\subsection{Phase return and representation of accumulated memory}
\label{vii:sec:retorno-memoria-gravedad}

The chronology of enriched routes distinguishes positional return,
orientational return, and phase return. The observable projection
identifies certain arrival states, while the cumulative record
preserves the sequence traversed. The geometric representation must
transport both data jointly.

\subsubsection{Inversion cocycle and change of sheet}
\label{vii:subsec:cociclo-memoria}

Let \(F_{\rm obs}\) be the chronological representation of TPK transport.
Its canonical blocks have the hierarchy
\[
\begin{array}{c|l}
27&\text{positional return with reversal of orientation},\\
54&\text{positional and orientational return},\\
108&\text{complete return of the observable phase}.
\end{array}
\]
The record counts reversals of orientation \(g(a)\) and
effective changes of sheet \(s(a)\). For a route,
\begin{equation}
 c(\gamma)=\sum_{a\subset\gamma}(g(a),s(a)),\qquad
 c(\gamma_2\gamma_1)=c(\gamma_2)+c(\gamma_1).
 \label{vii:eq:cociclo-ruta}
\end{equation}
The canonical computation of the block of length \(27\) yields
\(c_{27}=(1,18)\). Composition preserves those increments in
the four blocks completing the observable period.
On forward-oriented routes the record is cumulative;
its extension to the route groupoid uses \(\mathbb Z^2\) and
\(c(\bar\gamma)=-c(\gamma)\).

\begin{proposition}[Memory of the complete observable circuit]
\label{vii:prop:memoria-vuelta-completa}
Writing \(\mathcal K_{27}=F_{\rm obs}^{27}\), the lift
to the record, for \(x\in\mathcal Q_{\rm obs}\) and \(z\in\mathbb Z^2\), satisfies
\[
 \widetilde{\mathcal K}_{27}(x,z)
   =(\mathcal K_{27}x,z+(1,18)),\qquad
 \widetilde{\mathcal K}_{27}^{\,4}(x,z)
   =(x,z+(4,72)).
\]
After \(N\) complete circuits, the increment is \(N(4,72)\).
\end{proposition}

\begin{proof}
The observable projection of \(\mathcal K_{27}\) has order four.
The preceding canonical computation fixes \(c_{27}=(1,18)\) constantly
on that observable orbit. Additivity of
\eqref{vii:eq:cociclo-ruta} sums the four increments and gives
\(c_{108}=4(1,18)=(4,72)\).
Repetition applies the same concatenation law and produces
\(Nc_{108}\). Return of the first coordinate preserves the
increment of the second.
\end{proof}

\subsubsection{Affine representation of the cocycle}
\label{vii:subsec:representacion-afin-cociclo}

The discrete realization uses an element \(R_4\) of order four in
\(\operatorname{Spin}^+(1,3)\), a unit timelike vector \(u\)
fixed by its vector action, and a length unit \(\ell_0>0\).
The scale \(\ell_0\) preserves its provenance in the dimensional section.
We define
\begin{equation}
 \rho_{\rm C}^{\rm disc}(\gamma)=
 \bigl(R_4^{c_g(\gamma)},\,\ell_0c_s(\gamma)u\bigr)
 \in\operatorname{Spin}^+(1,3)\ltimes\mathbb R^{1,3}.
 \label{vii:eq:representacion-discreta-memoria}
\end{equation}
The power acts on the vector through the associated vector
representation of the spin group.

\begin{theorem}[Composition and translational amplitude of return]
\label{vii:thm:traslacion-memoria}
The map \(\rho_{\rm C}^{\rm disc}\) preserves the identity,
concatenation, and inversion. At complete returns,
\begin{equation}
 \rho_{\rm C}^{\rm disc}(\gamma_{108})=(I,72\ell_0u),\qquad
 \rho_{\rm C}^{\rm disc}(\gamma_{108}^{\,N})=(I,72N\ell_0u).
 \label{vii:eq:traslacion-retorno-108}
\end{equation}
Inversion reverses the sign of the translation.
\end{theorem}

\begin{proof}
The semidirect product is
\((R,a)(S,b)=(RS,a+Rb)\). Since \(R_4u=u\),
\[
\begin{split}
 \rho_{\rm C}^{\rm disc}(\gamma_2)
 \rho_{\rm C}^{\rm disc}(\gamma_1)
 &=
 \bigl(R_4^{c_g(\gamma_2)+c_g(\gamma_1)},
 \ell_0(c_s(\gamma_2)+c_s(\gamma_1))u\bigr)\\
 &=\rho_{\rm C}^{\rm disc}(\gamma_2\gamma_1).
\end{split}
\]
The cocycle of the identity is zero and that of the inverse is its negative.
For the complete circuit, substitute \(c_{108}=(4,72)\) and
\(R_4^4=I\). Concatenation of \(N\) circuits keeps the rotational
part equal to the identity and sums their translations.
\end{proof}

The amplitude \(72\ell_0\) is a length associated with a finite route.
A realization through a cotetrad and connection transports it to its
affine holonomy. The local contribution to the gravitational equations
is then obtained through the curvature of that connection,
variation of the material action, and normalization of its density.
The representation thus preserves the exact cumulative content
that must enter that realization.

% Article VII. Incorporation of pre-existing results, 10-09-2026.
% Complete source owners preserved in fuentes_conservadas:
% 17_pantalla_cubica_soldadura.tex (P), 18_clausura_energetica.tex (E).
% P: incidence, quotient, cellular soldering, and refinement.
% E: response, uniqueness, minimal extension, and tracial expectation.
% Previous dependencies of VII: APP--TRIT--TPK kernel; cubic realization
% of the route, transport cocycle, and compression q_vis=5/9.
% This fragment alone does not constitute the self-contained article.
% The screen identity is distinguished from the material spin contraction.

\subsection{Cubic incidence and boundary soldering}
\label{vii:sec:pantalla-respuesta}

The cubic realization of transport preserves six oriented incidences.
Its geometric quotient, Lorentzian form, and energetic response
are obtained from the same incidence matrix. Transport originates in
the enriched APP--TRIT--TPK route: the face, sign, carry, and
memory accompany each edge. This section develops the
linear operation acting on that boundary representation.

\subsubsection{Cubic incidence and the rank-four quotient}
\label{vii:sec:cociente-pantalla}

Let \(H_F=\mathbb R^F\), with
\(F=\{(i,\epsilon):1\leq i\leq3,\ \epsilon\in\{-1,1\}\}\),
its Euclidean inner product, and its basis \(e_{i,\epsilon}\).
Face opposition is the orthogonal involution
\(Re_{i,\epsilon}=e_{i,-\epsilon}\). Define
\[
 u=\sum_{i,\epsilon}e_{i,\epsilon},\qquad
 P_u=\frac{uu^{\mathsf T}}6,\qquad P_-=\frac{I-R}{2},\qquad
 P_0=\frac{I+R}{2}-P_u,\qquad P_\square=P_u+P_-.
\]
The vertices are \(V=\{-1,1\}^3\). The incidence matrix
\(M:\mathbb R^V\longrightarrow H_F\) is determined by
\[
 M_{(i,\epsilon),\nu}=\mathbf1_{\{\nu_i=\epsilon\}}.
\]

\begin{proposition}[Incidence decomposition]
\label{vii:prop:incidencia}
One has
\[
 MM^{\mathsf T}=12P_u+4P_-,
 \qquad
 \ker M^{\mathsf T}=P_0H_F.
\]
Consequently, the quotient
\(Q_\square=Q=H_F/P_0H_F\), orthogonally identified with
\(P_\square H_F\), has dimension four.
\end{proposition}

\begin{proof}
A face contains four vertices. Opposite faces have no common vertices,
and two faces on different axes have two common vertices.
Writing \(\mathbf J=uu^{\mathsf T}\) gives
\[
 MM^{\mathsf T}=4I+2(\mathbf J-I-R)
 =2(I-R)+2\mathbf J=4P_-+12P_u.
\]
The three projectors \(P_u,P_-,P_0\) are orthogonal and sum to \(I\);
their ranks are one, three, and two, respectively. The identity
\(\|M^{\mathsf T}x\|^2=\langle x,MM^{\mathsf T}x\rangle\)
identifies the stated kernel. The complementary rank is four.
\end{proof}

On \(Q\), the form
\[
 B=P_--P_u
\]
has signature \((3,1)\). Time orientation is fixed by \(u\).
With
\[
 t=\frac{u}{\sqrt6},\qquad
 d_i=\frac{e_{i,+}-e_{i,-}}{\sqrt2},\qquad
 W=\sqrt6P_u+\sqrt2P_-,
\]
one has \(We_{i,\epsilon}=t+\epsilon d_i\). These six vectors
are null for \(B\), whereas
\[
 3t+\nu_1d_1+\nu_2d_2+\nu_3d_3
\]
has squared norm \(-6\). The same incidence satisfies
\(MM^{\mathsf T}=2W^*W\) on \(Q\). The proper rotations of the cube
commute with the projectors, preserve both forms, and fix \(t\).
These identities determine the metric of the cubic representation.

\subsubsection{Cellular soldering and memory transport}
\label{vii:sec:soldadura-celular}

For each oriented edge \(a\), the reading of the route provides
the face \(\lambda(\underline a)\), the sign
\(\sigma(a)\), the source and target fibers, and the transport
\(U(a):Q_{s(a)}\longrightarrow Q_{t(a)}\).
Reversal satisfies \(U(\bar a)=U(a)^{-1}\).
Boundary half-edges preserve their formal reversal and their incidence
on the boundary. With \(\ell_m=\ell_0\,9^{-m}>0\), soldering is
\begin{equation}
 \beta_m(a)=\ell_m\sigma_m(a)
              n_{\lambda(\underline a)}^{s(a)},\qquad
 n_{i,\epsilon}=t+\epsilon d_i.
 \label{vii:eq:soldadura-arista}
\end{equation}
The face and sign are transported together with the route memory.
The reversal convention requires
\begin{equation}
 \beta_m(\bar a)=-U_m(a)\beta_m(a).
 \label{vii:eq:inversion-soldadura}
\end{equation}
Applying this relation twice recovers \(\beta_m(a)\);
the carry record separately preserves the ordered composition.

\begin{proposition}[Naturality of soldering under refinement]
\label{vii:prop:refinamiento-soldadura}
Let \(a=a_1\cdots a_9\) be a compatible refinement of an edge.
Suppose that the nine successor incidences, transported to the predecessor fiber,
preserve the face and sign, and that
\(\ell_{m+1}=\ell_m/9\). Let
\(\mathcal U_j:Q_{m,s(a)}\longrightarrow Q_{m+1,s(a_j)}\)
be the accumulated transport, including the refinement identification
and the path to the source of successor \(j\). Then
\[
 \beta_m(a)=
 \sum_{j=1}^{9}\mathcal U_j^{-1}\beta_{m+1}(a_j).
\]
The equality is compatible with reversal and with concatenation
of the memory record.
\end{proposition}

\begin{proof}
Each summand, expressed in the source fiber of the predecessor, is
\((\ell_m/9)\sigma_m(a)n_{\lambda(\underline a)}\).
The sum of nine terms reproduces the predecessor soldering.
Reversal reverses the order of the transports and replaces each
by its inverse; relation \eqref{vii:eq:inversion-soldadura} supplies
the overall sign. Memory is concatenated in the same order, so
refinement preserves both the vector representation and its
enriched record.
\end{proof}

% RESULTADO_RECUPERADO: S15, 11c_gravedad_holonomia_rev6.tex:128--212;
% S01, 02_dinamica_tres_hojas.tex:155--219; screen and response from30.
% Curvature, covariance, and Bianchi proofs assembled for autonomy.
% The smooth realization preserves its explicit holonomy conditions.

\subsection{Geometric realization of transport and soldering}
\label{vii:sec:realizacion-cartan}

The discrete representation of return and screen soldering
determine two aspects of transport: accumulation of memory along
routes and comparison of their fibres. Their geometric realization
preserves both operations. This section fixes the target
space, the dimensional normalization of the connection, and the identities
that its variation will use.

\subsubsection{Screen, cotetrad, and connection}

The screen quotient \(Q_\square\), its Lorentzian form, and the
soldering \(\beta_m\) come from
\ref{vii:sec:pantalla-respuesta}. In each cell, a screen
realization \(J_{{\rm scr},m}\) transports the form and soldering
to the corresponding geometric fibre. Its cellular image is
\[
 e_m=J_{{\rm scr},m}\beta_m.
\]
The scale \(\ell_m\) remains within the soldering. Frame
changes act jointly on \(J_{{\rm scr},m}\), the connection,
and cell transport.

A smooth realization of these data on an oriented manifold
\(\mathcal U\) of dimension four consists of a nondegenerate
cotetrad \(e=(e^I)\), a Lorentz connection
\(\omega=(\omega^I{}_J)\), and a route realization
\(\mathfrak R_{\rm C}\) satisfying
\begin{equation}
 \operatorname{Hol}_{\mathcal A_{\ell_0}}
       (\mathfrak R_{\rm C}(\gamma))
 =\rho_{\rm C}(\operatorname{Hol}_{\rm TPK}(\gamma)).
 \label{vii:eq:compatibilidad-holonomia-geometrica}
\end{equation}
In this matrix equality, the representation is normalized by
\[
 N_{\ell_0}(R,a)=(\Lambda(R),a/\ell_0),\qquad
 \rho_{\rm C}=N_{\ell_0}\circ\rho_{\rm C}^{\rm disc},
\]
where \(\Lambda:\operatorname{Spin}^+(3,1)\to SO^+(3,1)\)
is the vector representation. The map \(N_{\ell_0}\)
preserves the semidirect product: dividing the translation
\(a+\Lambda(R)b\) by \(\ell_0\) is equivalent to composing
the two normalized translations. The complete circuit therefore yields
\(72u\) in the normalized connection and \(72\ell_0u\) in the
dimensional chart. The spin lift preserves, separately,
\(R_4^{c_g(\gamma)}\); the vector representation has central
kernel \(\{1,-1\}\), so that lift is retained when
spinorial fields are used.
Identity, concatenation, inversion, and subdivision are maintained in the
corresponding realization. The screen pairing and holonomy condition
specify the data of the realization used in the field equations.

With \(\eta=\operatorname{diag}(-1,1,1,1)\), define
\(g=\eta_{IJ}e^I\otimes e^J\). For a positive elementary
scale \(\ell_0\), constant in the chart, the affine connection is
\begin{equation}
 \mathcal A_{\ell_0}
 =\begin{pmatrix}\omega&\ell_0^{-1}e\\0&0\end{pmatrix}.
 \label{vii:eq:conexion-afin-realizada}
\end{equation}
The cotetrad has dimensions of length; the factor
\(\ell_0^{-1}\) makes its translational component dimensionally homogeneous.

\begin{proposition}[Curvature and torsion of the realized connection]
\label{vii:prop:curvatura-afin}
The curvature of \eqref{vii:eq:conexion-afin-realizada} is
\begin{equation}
 \mathcal F_{\ell_0}
 =\mathrm d\mathcal A_{\ell_0}
   +\mathcal A_{\ell_0}\wedge\mathcal A_{\ell_0}
 =\begin{pmatrix}F_\omega&\ell_0^{-1}T\\0&0\end{pmatrix},
 \quad F_\omega=\mathrm d\omega+\omega\wedge\omega,
 \quad T=\mathrm de+\omega\wedge e.
 \label{vii:eq:curvatura-afin-realizada}
\end{equation}
\end{proposition}
\begin{proof}
The product of matrices of forms produces \(\omega\wedge\omega\)
in the diagonal block and
\(\ell_0^{-1}\omega\wedge e\) in the translational block.
The exterior derivative contributes
\(\mathrm d\omega\) and \(\ell_0^{-1}\mathrm de\), since
\(\ell_0\) is constant. Their sum proves the equality.
\end{proof}

Curvature describes the transport defect of frames and torsion
that of the cotetrad. Both belong to the same realized connection;
their components preserve distinct tensor domains.

\begin{proposition}[Covariance and Bianchi identities]
\label{vii:prop:bianchi-cartan}
Under a frame change \(g_L:\mathcal U\to SO^+(3,1)\),
\[
 e'=g_Le,\qquad
 \omega'=g_L\omega g_L^{-1}-\mathrm dg_L\,g_L^{-1},
\]
one has \(T'=g_LT\) and \(F_{\omega'}=g_LF_\omega g_L^{-1}\).
Moreover,
\begin{equation}
 D_\omega T=F_\omega\wedge e,\qquad D_\omega F_\omega=0.
 \label{vii:eq:bianchi-realizacion}
\end{equation}
\end{proposition}
\begin{proof}
In \(\mathrm d(g_Le)+\omega'\wedge g_Le\), the terms
containing \(\mathrm dg_L\wedge e\) cancel.
The same substitution into the curvature gives its transformation
by conjugation. For the first Bianchi identity, expand
\[
 D_\omega(\mathrm de+\omega\wedge e)
 =(\mathrm d\omega+\omega\wedge\omega)\wedge e.
\]
In \(\mathrm dF_\omega+\omega\wedge F_\omega-F_\omega\wedge\omega\),
the terms containing \(\mathrm d\omega\) and the cubic
products cancel in pairs. This proves the second identity.
\end{proof}

\subsubsection{The three quadratic realizations of curvature}

The TRIT sheet \(\tau\in\{+1,0,-1\}\) is represented in
the Cartan algebra by generators \(J_{ab}=-J_{ba}\)
and \(P_a^{(\tau)}\), with relations
\begin{align*}
 [J_{ab},J_{cd}]&=
 \eta_{bc}J_{ad}-\eta_{ac}J_{bd}
 -\eta_{bd}J_{ac}+\eta_{ad}J_{bc},\\
 [J_{ab},P_c^{(\tau)}]&=
 \eta_{bc}P_a^{(\tau)}-\eta_{ac}P_b^{(\tau)},\\
 [P_a^{(\tau)},P_b^{(\tau)}]&=-\tau J_{ab}.
\end{align*}
For constant \(\ell>0\), write
\[
 \mathcal A_\tau^{\rm Cart}
 =\tfrac12\omega^{ab}J_{ab}+\ell^{-1}e^aP_a^{(\tau)}.
\]

\begin{proposition}[Curvature of the trit-typed realization]
\label{vii:prop:curvatura-tritificada}
With the preceding relations,
\begin{equation}
 \mathcal F_\tau^{\rm Cart}
 =\tfrac12(F_\omega^{ab}-\tau\ell^{-2}e^a\wedge e^b)J_{ab}
     +\ell^{-1}T^aP_a^{(\tau)}.
 \label{vii:eq:curvatura-tritificada}
\end{equation}
\end{proposition}
\begin{proof}
Expand \(\mathrm d\mathcal A+\tfrac12[\mathcal A,\mathcal A]\).
The bracket of the Lorentz sector produces \(F_\omega\), the
mixed bracket produces \(D_\omega e\), and the translational bracket contributes
\(-\tfrac12\tau\ell^{-2}e^a\wedge e^bJ_{ab}\).
\end{proof}

The central sheet recovers the affine connection; the outer sheets
preserve the two signs of the constant-curvature term.
The representation thus maintains the route's tritic regime before
applying the variational equations. The holonomy equality
\eqref{vii:eq:compatibilidad-holonomia-geometrica} corresponds to the
central sheet. Each outer sheet uses the representation
\(\rho_{{\rm C},\tau}\) in its Cartan group and its respective
holonomy condition. The current and metric tensor use
the fields and representation of the same realization.


\subsection{Oriented transport of the connection and area}
\label{sec:ii-costura-barbero}
The parity of the incidence functional allows the charts
\(\gamma\) and \(-\gamma\) to be compared. In the canonical geometric realization,
let \(E^a_i\) be a densitized triad, \(\Gamma_a^i(E)\) its
intrinsic connection, and \(K_a^i\) the extrinsic curvature in the
chosen normal convention. The Ashtekar--Barbero connection is
\(\mathcal A_a^i=\Gamma_a^i(E)+\gamma K_a^i\).

\begin{proposition}[Involution between oriented charts]
\label{prop:ii-costura-barbero}
For \(\gamma\in\mathbb R\setminus\{0\}\), the map
\[
 \mathfrak b:(\gamma,K,E)\longmapsto(-\gamma,-K,E)
\]
is involutive. It preserves \(\mathcal A\) and reverses the form
\[
 \Omega_{KE}=\frac1\kappa\int_\Sigma
 \delta K_a^i\wedge\delta E^a_i,
\]
with \(\kappa\) fixed and the same boundary conditions.
\end{proposition}
\begin{proof}
Applying \(\mathfrak b\) twice restores all three coordinates.
Since \(E\) remains fixed, so does \(\Gamma(E)\).
The product of the two odd coordinates satisfies
\((-\gamma)(-K)=\gamma K\), preserving the connection.
In the space of variations, \(\delta(-K)=-\delta K\) and
\(\delta E\) remains fixed: the pullback of the form is
\(\mathfrak b^*\Omega_{KE}=-\Omega_{KE}\).
\end{proof}

The oriented area \(\mathcal A_{\rm or}=\gamma a_0N\) and its
physical magnitude \(\mathcal A_{\rm fis}=|\gamma|a_0N\), for
\(a_0>0\) and \(N\geq0\), agree in the positive chart.
Upon prolongation to the opposite chart, the former is odd and the latter
remains invariant. This distinction preserves both the orientation
of incidence and positivity of the geometric observable.
The proposition establishes algebraic transport of these charts;
its realization as reversal of the spatial normal additionally uses
transport of matter and of the boundary data
of action~\eqref{eq:holst}.
